\documentclass{book}
\usepackage{amsmath}
\usepackage{amssymb}
\usepackage{amsthm}
\usepackage{hyperref}

\theoremstyle{definition}
\newtheorem{definition}{Definition}[section]
\newtheorem{proposition}[definition]{Proposition}

\setcounter{chapter}{1}
\setcounter{section}{4}

\begin{document}

\section{Dimension}

\begin{definition}
Let $V$ be a vector space over a field $F$. A \textbf{subspace chain} in $V$ is a finite sequence of subspaces
\[
\{0\} = V_{0} \subsetneq V_{1} \subsetneq \dots \subsetneq V_{k} \subseteq V.
\]
\end{definition}

\begin{definition}
For a subspace $U \subseteq V$, the \textbf{dimension} of $U$ is
\[
\dim U = \sup\,\{\,k \mid \{0\} = V_{0} \subsetneq \dots \subsetneq V_{k} = U \text{ is a subspace chain}\,\}.
\]
If the supremum does not exist, $\dim U = \infty$.
\end{definition}

\begin{proposition}
For any subspace $U \subseteq V$,
\[
\dim U + \dim(V / U) = \dim V.
\]
(Here $\infty + n = n + \infty = \infty + \infty = \infty$.)
\end{proposition}

\begin{proof}
Let $\pi \colon V \to V / U$ be the quotient map.

\textbf{$\dim V \le \dim U + \dim(V / U)$.}
Take any subspace chain $\{0\} = V_{0} \subsetneq \dots \subsetneq V_{k} = V$ in $V$.
For each $i$, either $V_{i-1} \cap U \subsetneq V_{i} \cap U$, or
$\pi(V_{i-1}) \subsetneq \pi(V_{i})$ (or both): if neither were strict, then
$V_{i-1} \cap U = V_{i} \cap U$ and $V_{i-1} + U = V_{i} + U$; picking
$x \in V_{i} \setminus V_{i-1}$ and writing $x = y + u$ with $y \in V_{i-1}$,
$u \in U$ (from the second equality) gives $u = x - y \in V_{i} \cap U
= V_{i-1} \cap U \subseteq V_{i-1}$, so $x = y + u \in V_{i-1}$, contradiction.

Thus the $k$ steps are split into those where the intersection with $U$
strictly grows, and those where the image in $V / U$ strictly grows.
After removing redundancies, the intersections form a chain in $U$,
the images a chain in $V / U$, so $k \le \dim U + \dim(V / U)$.
Taking the supremum over all chains in $V$ yields the inequality.

\textbf{$\dim V \ge \dim U + \dim(V / U)$.}
Take chains $\{0\} = U_{0} \subsetneq \dots \subsetneq U_{a} = U$ in $U$ and
$\{0\} = W_{0} \subsetneq \dots \subsetneq W_{b} = V / U$ in $V / U$.
Let $V_{j} = \pi^{-1}(W_{j})$; then $U = V_{0} \subsetneq \dots \subsetneq V_{b} = V$.
Concatenating,
\[
\{0\} = U_{0} \subsetneq \dots \subsetneq U_{a} = U
      = V_{0} \subsetneq \dots \subsetneq V_{b} = V
\]
is a chain in $V$ with $a + b$ strict inclusions.
Thus $\dim V \ge a + b$ for any such $a$, $b$; taking suprema gives
$\dim V \ge \dim U + \dim(V / U)$.
\end{proof}

\begin{proposition}
For any linear map $f \colon V \to W$,
\[
\dim(\ker f) + \dim(\operatorname{im} f) = \dim V.
\]
\end{proposition}

\begin{proof}
By Proposition 1.5.3 with $U = \ker f$,
$\dim(\ker f) + \dim(V / \ker f) = \dim V$.
By \href{s04.html}{Proposition 1.4.5}, $V / \ker f \cong \operatorname{im} f$ via an isomorphism
$\Phi$. Applying $\Phi$ (resp.\ $\Phi^{-1}$) to each term of a chain in $V / \ker f$
(resp.\ $\operatorname{im} f$) gives a chain of equal length in $\operatorname{im} f$
(resp.\ $V / \ker f$); hence $\dim(V / \ker f) = \dim(\operatorname{im} f)$.
Substituting yields the result.
\end{proof}

\end{document}
